%\usepackage{graphicx}
\chapter {Introduction}
\pagenumbering{arabic}

Computer research nowadays aimed at facilitating computer human interaction in every way. Pen-based interfaces give the user a pencil-paper like feeling that enhances interaction more than the current used keyboard-mouse computer interfaces. However, until now there are hardly any complete pen based computer systems. Generally, Current Pen-Based systems make use of the Pen or stylus to perform the same role as the mouse / keyboard. A minority of such systems are equipped with handwriting or digit recognition that is typically restricted for specific language. 

\section{Sketch Recognition}
Sketch recognition is defined as the process of identifying symbols that users draw using single or multiple strokes. Users draw strokes using a pen and the system should immediately interprets their strokes into objects that can be easily manipulated.

Sketchs here stand for any hand drawn drawing the user draws using digital stylus or similar device. Normally, sketch are composed of smaller segments we call strokes. The time between pen-up and pen-down events identify each stroke.  The strokes are mainly the path of points, which the user path through between the pen-down and pen-up events while using the stylus or similar device. 

 
%Sketch recognition is procedure that will input a sequence of points from user then segment them into clusters of strokes and recognize those clusters to primitive geometric objects and shapes. However, Sketch understanding is transforming the sketch from a group of primitive shapes into comprehendible identified symbols consistent with each other.
% the entire sketched symbols.
%Sketch recognition is defined as the process of identifying symbols that users draw using single or multiple strokes. Users draw strokes using a pen and the system should immediately interprets their strokes into objects that can be easily manipulated. This paper uses Particle Swarm Algorithm (PSO) to divide the strokes the user draws into meaningful geometric primitives. These geometric primitives are grouped to formulate symbols which are further identified. The results show that using PSO improves segmentation results which guide the symbol recognition phase. This paper uses Support Vector Machine(SVM) classifier which further improves the final recognition accuracy.  %u Final recognition has improved accuracy after the PSO algorithm using a SVM classifier.  
 %Hand drawn sketches are the natural method of exchanging thoughts between engineers and humans in general. Nevertheless, researchers barely give attention to improve current gesture and sketch recognition systems. In a sketch understanding system, users need to feel the freedom of paper and pencil along with the flexibility of computer manipulation. 
 
 %A picture speak a thousand words is also true for sketches. 
%Scientists generally and engineers specifically express thoughts and designs using sketches. Engineers use sketches to exchange designs as a natural method of communication rather than writing or speaking. Engineers use paper and pencil design in early stages of design. Design engineers need a powerful computer based system with symbol recognition to help them design, manipulate and store sketches more effectively than using only papers. Increasing the interaction between computers and users in sketch and CAD systems has been the reason for the emerging of few advanced sketch recognition systems \cite{CAD}. Sketch recognition system converts the drawings and informal graphs user draws into objects and symbols that can be used in further interaction with the system. 

%Sketch recognition is divided into four main steps preprocessing, segmentation, symbol recognition and finally sketch understanding. The preprocessing process saves user input stroke points from the hardware and collects basic information about the stroke then proceeds to remove noise and primary data calculation. In segmentation phase, the system divides the stroke into a set of simple geometrical primitives or segments. The third phase, sketch recognition, is clustering strokes and segments to formulate symbols that can be recognized by a classifier system. Finally, the sketch understanding phase uses priori information about the sketch to support identification of symbols by excluding any possible invalid symbols. 
%the basic symbols are identified in consistent with the whole sketch. \\
%from the user information from the users converts the strokes the user draw into segments which is recognized into basic symbols that can be easily manipulated. Finally, the basic symbols are interpreted in consistent with the whole sketch. \\
%This paper introduces a sketch recognition system. The presented system uses' particle swarm algorithm (PSO) to optimally segment users strokes. Users can draw symbols using any number of strokes in any order they like. The system segments stroke user draws then cluster symbols together then passes them to the SVM symbol recognizer. The experiments in section \ref{sec:Experiments} shows that the proposed system improves the final symbol recognition accuracy. 

%The remaining of this papers is organized as follows; section \ref{Sec:RelatedWork} provides a brief literature review.  Section \ref{PSO} explains the general particle swarm algorithm. The proposed system block diagram along with a detailed system description is presented in Section \ref{Sysdisc}. The experiments and results achieved are explained in Section \ref{sec:Experiments}. Finally, section \ref{ConclusionandFutureWork} states the conclusion and future work. 
%The next section will relate previous work in the field. The remaining of the paper will be as following; Section \ref{PSO} will explain the general particle swarm algorithm which is used in the presented system. Section \ref{Sysdisc} will describe the system block diagram and then section \ref{Prepross} will explain the preprocessing step, section \ref{seg} will describe the segmentation algorithm in details. Finally, section \ref{sec:Experiments} will describe the experiments that was preformed.\\
 
  
 
\section{Importance of sketch recognition}
\label{sec:ImportanceOfSketchRecognition}


Scientists generally and engineers specifically express thoughts and designs using sketches. Engineers use sketches to exchange designs as a natural method of communication rather than writing or speaking. Engineers use paper and pencil design in early stages of design. The sketching process itself go thorough numerous stages. Firstly, early stages with rough drafting designs or notations which are mainly done using paper and pencil. After those primary rough sketches are edited and altered until they are stabilized and finally they are transformed into the computer. 

%Design engineers need a powerful computer based system with symbol recognition to help them design, manipulate and store sketches more effectively than using only papers. Increasing the interaction between computers and users in sketch and CAD systems has been the reason for the emerging of few advanced sketch recognition systems \cite{CAD}. Sketch recognition system converts the drawings and informal graphs user draws into objects and symbols that can be used in further interaction with the system. 

%People normally sketch thoughts rather than expressing them linguistically. For instance engineers mostly drafts their early designs using paper and pencil before using CAD tools.  People in general and engineers specially make use of various sketches during design process as instance organizational charts, informal drafting and formal design sketches \cite{}. The sketching process itself go thorough numerous stages. Firstly, early stages with rough drafting designs or notations which are mainly done using paper and pencil. After those primary rough sketches are edited and altered until they are stabilized and finally they are transformed into the computer.


 Paper and Pencil are preferred in design as they offer a remarkable ease in editing and altering the primary rough sketches. Current CAD systems do not give the designer the ease of altering the early designs. Designer had to sketch primary designs on paper and pencil next, the finally design is transferred into computer using a menu CAD system. The process is time consuming and tiring. The need for fast design and development process is driven by the fast growing industries and technology.
 
 A sketch recognition and understanding system will provide a powerful tool to the designers to draw sketch as they usually draw using paper and pencil with the advantages of computer program. Designer needs a normal sketching environment that will provide him facilities in paper interface and the modifying and simulation that is in current CAD systems. 
 
\section{Applicability }
% talk about application like cad systems , software engineering uml diagrams ....
% meeting idea brainstorming , mind mapping string ..
% communication program  
% sketching programs ...
Application of sketch recognition are ranging from engineering device simulation and design. Drawing and sketching programs are most likely to benefit from integrating a sketch recognition system into there design environments. Other applications include object retrieval, map navigations and various engineering design systems.

 
\section{Research Scope}
%rewirte this part.............................
%  research will invistgated till sketch recognition in domain of uml diagrams. 
%This research will include investigating the best method to recognize sketches.  The scope of the research will include recognizing various sketches. The methods used in this research are comparing current systems and enhancing them and adobpiting them into simple geometrical symbols.  Using well established AI methods specially \textit{Swarm Intelligence} to solve the problem. 
 



\section{Problem Definition}
\label{sec:ProblemsAndChallenges}


%Sketch recognition is divided into three main steps preprocessing, segmentation and symbol recognition. The preprocessing phase captures user input stroke points and collects basic information about the stroke then proceeds to remove noise and compute basic statistical and geometrical information. In segmentation phase, the strokes are divided into a set of simple geometrical primitives or segments. In the third phase, sketch recognition, strokes and segments are clustered to formulate symbols that can be recognized by a classifier system.


 Understanding hand drawn sketches may be trivial for humans but it is challenging problem when speaking about computer-based system. Most of the research made on the sketch understanding is achieved for one particular domain or for strictly a set of pre defined symbols. Further, others restrict the user the freedom in drawing which loose the main goals of paper-like-interface systems.  The only few systems that have achieved the desired degree of freedom for the users are extremely computationally expensive.   

The research in the area had faced few problems that will need to be handled to implement an entirely sketch understanding system. The user draws strokes in a spontaneous and ambiguous manner that intensifies the challenges. The first problem is segmenting strokes from the way the user draw it into primitive geometric. Ink parsing is the problem dealing with clustering interpolated strokes and grouping them into symbols. After that, recognizing the symbols from the space the system can identify is the final problem. All those problems must be solved on a real time constrain to maintain an online sketching system if it was a requirement.

The next sections will explain and investigate each of those problems in details.




\subsection{Segmentations}
\label{sec:Segmentations}

Segmentation is dividing the strokes draws into the geometric primitive's as lines, arcs, and ellipses. The strokes as mentioned before are the path of points between pen-down and pen-up events while using a stylus. The processing of those points is done after that to firstly to detect the critical points or vertices. After that the segment is classified as either line, arc, ellipse or circle. The main problem in this process is in detecting the critical and vertices which are points divide strokes into lines and arc. The user draw the lines and arc in an undefined manner which make it difficult for the program to deicide where the user had finish a line and started a new arc or vice verse; Fig [\ref {fig:strokeinterpolation}].


\begin{figure}

		\includegraphics[scale=0.8]{../../neededfiles/Figures/strokeinterpolation.eps}
	\caption[Segmentation error]{The figure shows an arc that the user draw which can be segmented in different ways as a single line or as an arc where a) The stroke users draw  b)interpretation as a single line segment  b) interpretation as an arc.}
	\label{fig:strokeinterpolation}
\end{figure}



		

\subsection{Ink parsing}
\label{sec:InkParsing}

 The clustering of segmented strokes into group of strokes before passing into the next stage this process is called ink parsing. It is simply deciding when a symbols ends and a new symbol begins.  Ink parsing is most important and difficult challenge as grouping of strokes can be done in many ways and configuration, which increase exponentially with the number of strokes. In addition, the ink parsing will affect the next steps of recognizing. The user does not draw the symbol in same order. Moreover, users can draw a sequence of strokes in symbols then leave it to draw another then return again to add segment to symbols so temporal data is helpful but not conclusive.  The user may also draw the same symbols using different orders and methods for example the user can draw a rectangle using single stroke or using two, three or even four strokes. All those factors gather with the sloppy drawing of the user to increase the complexity of the problem.

\subsection{Ambiguities}
\label{sec:Ambiguities}
 	
The major problem of sketch understanding is the high level of ambiguities in the drawing. The user draws strokes that the computer cannot decide if it is arc or line. The nature of hand drawn sketch it self is highly ambiguous even for human eye you may see two drawn you cannot decide if it is square or rectangle.
 The ambiguities in sketch understanding are found at many levels in the system. Firstly, at the segmentation phase where the same user stroke can be interpreted as a ellipse and a circle. Then the symbol recognition as a symbol can have more than one interpolation. Due to the user sloppy drawing the symbols can have more than one interpolation which varies with the domain drawn in it similar problem was faced in handwriting recognition see Fig [\ref{fig:VorU}] but it was handled with knowing information about the context information which may not be the case in sketch recognition. 
 
\begin{figure}	
	\centering
		\includegraphics[scale=0.85]{../../neededfiles/Figures/VorU.eps}
	\caption[Ambiguities Handwriting ]{The figure shows that the same letter can be interpreted differently in various contexts; first as V in drive or as U in Cut}
\label{fig:VorU}
\end{figure}

The ambiguity problem must be first recognized ambiguities or varicose interpolations of same symbol. After that, resolve those ambiguities is detecting the best correct interpolation of the symbol.  To solve this problem the program must be able to detect the amount of information needed to decide which interpolation of symbols is the best one. 


\section{Thesis Overview}
This thesis is arranged such that chapter \ref{} presents a survey of different components of the
sketch recognition. In chapter \ref{}, more light is shed on the feature filter technique for dimensionality
reduction. This is in conjunction with a detailed description of the proposed methods to
enhance the feature filtering. Comparative studies among the proposed methods and stateof-
arts methods are presented in chapter 4. Finally, chapter 5 presents the conclusions of
this thesis and suggested future work.
%\subsubsection{Search space}
%\label{sec:Searchspace}
%
%
%The user drawn symbols can be having more than one interpolation. In addition, those interpolations increase exponentially with the number of strokes the user draw and as the sketch expands the number of interpolation escalate. This intensifies the search space that the program search to resolve the ambiguities and find best interpolation.
%
%\sublsubsection{Real time}
%\label{sec:Realtime}
%
%
%There are mainly three types of sketch understanding system offline \cite{Vibratory8}, online and interactive. Both online and interactive recognition like \cite{incrmentintention41, SmartSketch56} need the processing after each stroke to be real time so the user does not feel the processing.  



%A picture speak a thousand words is also true for sketches. 
%Scientists generally and engineers specifically express thoughts and designs using sketches. Engineers use sketches to exchange designs as a natural method of communication rather than writing or speaking. Engineers use paper and pencil design in early stages of design. Design engineers need a powerful computer based system with symbol recognition to help them design, manipulate and store sketches more effectively than using only papers. Increasing the interaction between computers and users in sketch and CAD systems has been the reason for the emerging of few advanced sketch recognition systems \cite{CAD}. Sketch recognition system converts the drawings and informal graphs user draws into objects and symbols that can be used in further interaction with the system. 

%Sketch recognition is divided into four main steps preprocessing, segmentation, symbol recognition and finally sketch understanding. The preprocessing process saves user input stroke points from the hardware and collects basic information about the stroke then proceeds to remove noise and primary data calculation. In segmentation phase, the system divides the stroke into a set of simple geometrical primitives or segments. The third phase, sketch recognition, is clustering strokes and segments to formulate symbols that can be recognized by a classifier system. Finally, the sketch understanding phase uses priori information about the sketch to support identification of symbols by excluding any possible invalid symbols. 
%the basic symbols are identified in consistent with the whole sketch. \\
%from the user information from the users converts the strokes the user draw into segments which is recognized into basic symbols that can be easily manipulated. Finally, the basic symbols are interpreted in consistent with the whole sketch. \\